\homework{3}{Fri 22 Sep 2023 23:18}{}

\begin{problem}
	Consider the initial value problem
\[
u_t=\sin u_x ; \quad u(x, 0)=\frac{\pi}{4} x .
\]
Verify that the assumptions of the Cauchy-Kovalevskaya theorem are satisfied and obtain the Taylor series of the solution about the origin.
\end{problem}
\begin{proof}
	Differentiate, we have
	\[
		u_{tt} = \cos (u_x) u_{x t}
	.\] 
	Now $\cos$ is an analytic function, so the assumptions are satisfied. To compute derivatives at $(0,0)$, 
	\[
		u_x(x,0) = \frac{\pi}{4}	.\] 
	Now
	\[
		u_{t x} = \cos(u_x) u_{x x} = 0
	.\] 
	This implies $u_{t t} (0,0) = 0$.
	So
\[
D^{(\alpha_1, \alpha_2)} u (0,0) = 0 \text{ for }\alpha_1 > 1 \text{ or } \alpha_2 > 1
.\] 
Hence the taylor series of the solution about the origin is
\[
	u(x,t) = \frac{\pi}{4} x + \frac{1}{\sqrt{2} } t
.\] 

\end{proof}



\begin{problem}
	Consider the Cauchy problem for $u(x, y)$
\[
\begin{aligned}
& u_y=a(x, y, u) u_x+b(x, y, u) \\
& u(x, 0)=0
\end{aligned}
\]
Let $a$ and $b$ be analytic functions of their arguments. Assume that $D^\alpha a(0,0,0) \geq 0$ and $D^\alpha b(0,0,0) \geq 0$ for all $\alpha$. (Remember by definition, if $\alpha=0$ then $D^\alpha f=f$.)\\
(a) Show that $D^\beta u(0,0) \geq 0$ for all $|\beta| \leq 2$.\\
(b) Prove that $D^\beta u(0,0) \geq 0$ for all $\beta=\left(\beta_1, \beta_2\right)$. (Hint. Argue as in the proof of the Cauchy-Kovalevskaya theorem; i.e., use induction in $\beta_2$ )
\end{problem}
\begin{proof}
	$u(0,0) = 0$ is given. We have $u_x = u_{x x} = 0$ at $(0,0)$. To compute $u_{yx }$ and $u_{yy }$, differentiate the PDE:
	\begin{align*}
		u_{yx }(0,0) &= D_x a \, u_x + D_u a \, u_x^2 + a u_{x x} + D_x b + D_u b \, u_x\\
		&= D_x b \ge  0
	.\end{align*}
	\begin{align*}
		u_{yy }(0,0) &= D_y a u_y + D_u a \, u_x u_y + a u_{x y} + D_y b + D_u b \, u_y \\
		&= D_y a \, b + D_x b + D_y b + D_u b \, b 
		\ge 0
	.\end{align*}

	No we prove the second part. For $\beta_2 = 0$, we have $D^\beta u(0,0) = \frac{\partial ^{\beta_1}}{\partial x^{\beta_1}} u = 0$. This proves the base case. Now assume inductively that $D^\beta u(0,0) \ge 0$ for all $\beta$ such that $\beta_2 \le  k$. Now set $\beta_2 = k$ and we have
	\begin{align*}
		D^{(\beta_1, k_1)} u 
		&= D^{\beta} u_y \\
		&= \sum_{\gamma \le \beta} \binom{\beta}{\gamma} D^\gamma a D^{\beta - \gamma} u_x + D^\beta b\\
		&= \sum_{\gamma \le \beta} \binom{\beta}{\gamma} D^\gamma a D^{\beta - \gamma + (1,0)}u + D^\beta b
	.\end{align*}
	Now $D^\gamma a$ and $D^\beta b \ge 0$  by assumption, and $D^{\beta - \gamma + (1,)}$ has no more than $k$ derivatives on $y$, so the induction hypothesis gives  $D^{\beta - \gamma + (1,0)}u\ge 0$. The induction is complete and the claim is proved.
\end{proof}
\newpage

\begin{problem}
	(Kovalevskaya's example) Show that the line $\{t=0\}$ is characteristic for the heat equation $u_t=u_{x x}$. Show there does not exist an analytic solution $u$ of the heat equation in $\mathbb{R} \times \mathbb{R}$, with $u=\frac{1}{1+x^2}$ on $\{t=0\}$. (Hint: Assume there is an analytic solution, compute its coefficients, and show that the resulting power series diverges in any neighborhood of $(0,0)$.)
\end{problem}
\begin{proof}
	First compute the derivatives of $\frac{1}{1+x^2}$:
	\[
		\frac{d ^n}{d x^n} \frac{1}{1+x^2}
		= \frac{n!}{(x^2 + 1)^{n+1}} 
		\sum_{k=0}^{\left\lfloor n / 2 \right\rfloor}  \binom{n+1}{2k+1} x^{n - 2 k - 1}(-1)^{n+k}
	.\] 

	Now observe that
	\[
		D^{(\beta_1, \beta_2)} u = D^{(\beta_1, \beta_2 - 1)} u_t =D^{(\beta_1, \beta_2 - 1)} u_{x x} = D^{(\beta_1+2, \beta_2-1)} u
	.\] 
	In particular
	\[
		D^{(0,\beta_2)} u = D^{(2 \beta_2, 0)} u
	.\] 
	Now assume $u$ is analytic at $(0,0)$ so that we have the power series expansion
	\[
		u(x, t) = \sum_{|\beta|=0}^{\infty} D^{\beta} u \frac{1}{\beta!}x^{\beta_1} t^{\beta_2}
	.\] 
	We restrict ourselves to the line $x = 0$. So
	\[
		u(0,t) = \sum_{n = 0}^{\infty } \frac{1}{n!} D^{(2 n, 0)}u(0,0)t^n
	.\] 
	Where
	\[
		D^{(2n, 0)}u(0,0) = (2n)!
	.\] 
	Compute the radius of convergence:
	\[
		\frac{1}{R} = \limsup_{n \to \infty } \left( \frac{(2n)!}{n!} \right) ^{\frac{1}{n}} \ge \limsup_{n \to \infty } n = \infty 
	.\] 
	Hence $R = 0$, i.e. this power series is divergent in any neighborhood of $(0,0)$.


\end{proof}



